\textbf{Proof.}
Because \(k\le d\le n^2\),
\[
 0<\alpha^2=\frac{\sqrt{k}}{16n}\le\frac1{16},
\]
so \(0<\alpha\le1/4\).  Hence all coordinates of every \(\mathbf P_\sigma\) are nonnegative.  Each displayed pair sums to \(2/k\), and there are \(k/2\) pairs, so \(\mathbf P_\sigma\in\Delta_d\).  Direct calculation gives
\[
 \lVert\mathbf P_\sigma-\mathbf Q\rVert_1
 =\sum_{j=1}^{k/2}
   \left(\frac{\alpha}{k}+\frac{\alpha}{k}\right)
 =\alpha.
\]

Let \(M=2^{-k/2}\sum_\sigma\mathbf P_\sigma^n\), a probability law on \(\{1,\ldots,k\}^n\).  For independent uniform sign vectors \(\sigma\) and \(\tau\), expansion of the squared mixture likelihood ratio under \(\mathbf Q^n\), followed by independence of the \(n\) coordinates, gives
\[
 \begin{aligned}
 1+\chi^2(M,\mathbf Q^n)
 &=E_{\sigma,\tau}\left[
   \sum_{z=1}^{k}\frac{P_{\sigma,z}P_{\tau,z}}{Q_z}
  \right]^n\\
 &=E_{\sigma,\tau}\left[
  1+\frac{2\alpha^2}{k}
     \sum_{j=1}^{k/2}\sigma_j\tau_j
  \right]^n.
 \end{aligned}
\]
This identity is well defined because all laws have common support \(\{1,\ldots,k\}\).  The bracket is at least \(1-\alpha^2>0\).  The elementary inequalities \(1+t\le e^t\) for \(t>-1\) and \(\cosh t\le e^{t^2/2}\), together with independence of the Rademacher products \(\sigma_j\tau_j\), imply
\[
 \begin{aligned}
 1+\chi^2(M,\mathbf Q^n)
 &\le
 \left\{\cosh\left(\frac{2n\alpha^2}{k}\right)\right\}^{k/2}\\
 &\le \exp\left(\frac{n^2\alpha^4}{k}\right)
 =e^{1/256}.
 \end{aligned}
\]
For completeness, the second elementary inequality follows by comparing the power series of \(\cosh t\) and \(e^{t^2/2}\).  Cauchy--Schwarz applied to the likelihood ratio therefore yields
\[
 \lVert M-\mathbf Q^n\rVert_{\mathrm{TV}}
 \le\frac12\sqrt{\chi^2(M,\mathbf Q^n)}<\frac12.
\]

We next convert indistinguishability into squared-error risk.  The second multinomial stream has law \(\mathbf Q^n\) under both the null pair \((\mathbf Q,\mathbf Q)\) and every alternative pair \((\mathbf P_\sigma,\mathbf Q)\).  Tensoring with this common law leaves total variation unchanged.  For an arbitrary estimator \(\widehat L\) based on both streams, define \(\varphi=\mathbf 1\{\widehat L\ge\alpha/2\}\).  The target is zero at the null and equals \(\alpha\) at every alternative.  If \(R(\widehat L)\) is the maximum squared risk of \(\widehat L\), then
\[
 \begin{aligned}
 2R(\widehat L)
 &\ge E_{\mathbf Q^n\otimes\mathbf Q^n}(\widehat L^2)
   +E_{M\otimes\mathbf Q^n}\{(\widehat L-\alpha)^2\}\\
 &\ge\frac{\alpha^2}{4}
 \left[
  P_{\mathbf Q^n\otimes\mathbf Q^n}(\varphi=1)
  +P_{M\otimes\mathbf Q^n}(\varphi=0)
 \right]\\
 &\ge\frac{\alpha^2}{4}
 \left\{1-\lVert M-\mathbf Q^n\rVert_{\mathrm{TV}}\right\}
 \ge\frac{\alpha^2}{8}.
 \end{aligned}
\]
Taking the infimum over \(\widehat L\) proves
\[
 \mathfrak R^{L_1}_{n,d}\ge\frac{\alpha^2}{16}
 =\frac{\sqrt{k}}{256n}.
\]
The exact fixed-sample Markov reduction in \(\mathrm{lem:l1-lower-transfer}\) gives
\[
 \mathfrak R_{n,d,\epsilon}
 \ge\frac1{16}\mathfrak R^{L_1}_{n,d}.
\]
Finally, \(k=2\lfloor d/2\rfloor\ge d/2\) for \(d\ge2\), whence
\[
 \mathfrak R_{n,d,\epsilon}
 \ge\frac{\sqrt{k}}{4096n}
 \ge\frac{\sqrt d}{4096\sqrt2\,n}.
\]
The embedded causal laws have propensity \(1/2\), so \(\mathrm{ass:fixed-overlap}\) holds for every \(0<\epsilon<1/2\); their causal samples are i.i.d. by the reduction.  This verifies the support, positivity, overlap, and sample-size conditions used above.